\documentclass[11pt]{article}
\usepackage[margin=1in]{geometry}
\usepackage{amsmath,amssymb,amsthm,mathtools}
\usepackage[hidelinks]{hyperref}
\newtheorem{theorem}{Theorem}
\newtheorem{lemma}[theorem]{Lemma}
\newtheorem{corollary}[theorem]{Corollary}
\newcommand{\Q}{\mathbb Q}
\newcommand{\Qb}{\overline{\mathbb Q}}
\newcommand{\C}{\mathbb C}
\newcommand{\Lfac}{\mathcal L}
\title{Two angular coordinates: a rank-two comparison and integer-Gamma families}
\author{Internal proof and exact certificates; no priority claim}
\date{16 November 2026}
\begin{document}
\maketitle
\noindent\textbf{Result and provenance.}
The pencil theorem in \href{SIMPLEX_PENCILS_AND_AFFINE_LIMITS.tex}{the preceding note}
does not cover a general phase in two angular coordinates. Here we prove a
specific two-coordinate comparison directly, and use a good-reduction test
to obtain a degree-four factorial/Gamma family in every dimension, for
every positive integer shape vector. All sufficiently large moments survive
at the stated algebraic scales, despite arbitrary complex lower terms.

The comparison mechanism is credited to Long--Le Blanc,
\href{https://raw.githubusercontent.com/long-mathematics/exponential-integral-theorem/92db79eedeb064694ff20e9201d73fc2bc0857f4/papers/05-polynomial-flags/polynomial-flag-exponential-periods.tex}{EIT V, linear criterion and strict-extremum lemma},
at commit \texttt{92db79eedeb064694ff20e9201d73fc2bc0857f4}.
Their \href{sources/long-eit-vi-current.tex}{EIT VI} already proves odd-degree
and quartic factorial families, including a different all-dimensional
characteristic-five argument. We do not claim these general ideas or the
original families as new. Our cubic calculation replaces V's geometric
irreducibility input by an elementary rational Riccati obstruction and
uses reduction at $h=0$ for all three integer shape exponents at once.
Thus this note does not depend on the unreviewed stronger Galois,
boundary-field independence, saddle-continuation or multicolumn claims.
The arithmetic external input is Beukers' ordinary-point linear lifting,
already checked in \href{LONG_EIT_AND_FC2_AUDIT.tex}{the EIT I audit}.
This is an internal written proof, not independent specialist review or
formal verification. No novelty claim is made.

\section{A free polynomial twisted quotient}
Set
\[
 q(X,Y)=X^3-3X+Y^2,\qquad
 T=\{0\le X\le2,\ |Y|\le1-X/2\},\qquad h=1/z.
\]
Use the additive quotient by
\[
 D_X=h\partial_X+3(X^2-1),\qquad D_Y=h\partial_Y+2Y.
\]
This is a quotient of vector spaces, not a quotient algebra. For every
$r\in\Qb[X,Y]$ there are unique $a_r,b_r\in\Qb[h]$ and some polynomial
$P_r,Q_r\in\Qb[h,X,Y]$ such that
\[
 r=a_r(h)+b_r(h)X+D_XP_r+D_YQ_r.                 \tag{1}
\]
For rational amplitudes with denominators supported at $2,3$, all these
polynomials belong to $\mathbb Z[1/6,h,X,Y]$. Indeed the recursions are
\begin{align*}
 N(X^u)&=N(X^{u-2})-\frac{h(u-2)}3N(X^{u-3})&& (u\ge2),\\
 N(Y^v)&=-\frac{h(v-1)}2N(Y^{v-2})&& (v\ge1),
\end{align*}
with $N(1)=1$, $N(X)=X$ and zero-coefficient terms omitted. They follow by
subtracting $D_X(X^{u-2}/3)$ or $D_Y(Y^{v-1}/2)$. The two one-variable
reductions commute. Their tensor product annihilates $D_XP+D_YQ$ and fixes
$1,X$, proving uniqueness, including after every specialization $h=1/\xi$
with $\xi\ne0$. At $h=0$ it is precisely reduction modulo
$(X^2-1,Y)$.

Writing $U_0(z)=\int_Te^{zq}\,dX\,dY$ and
$U_1(z)=\int_TXe^{zq}\,dX\,dY$, Stokes gives
\[
 \int_T r e^{zq}=a_r(1/z)U_0+b_r(1/z)U_1+
 h\int_{\partial T}e^{zq}(P_r\,dY-Q_r\,dX).       \tag{2}
\]
No spatial denominator or extra boundary is introduced.

\section{Elementary irreducibility and actual comparison}
Direct instances of (1) give
\begin{align*}
 q&=-5h/6-2X+D_X(X/3)+D_Y(Y/2),\\
 qX&=-2-7hX/6+D_X((X^2-2)/3)+D_Y(XY/2).
\end{align*}
Consequently the actual column $U=(U_0,U_1)^t$ satisfies
\[
 U'=A(z)U+\text{actual boundary forcing},\qquad
 A(z)=\begin{pmatrix}-5/(6z)&-2\\-2&-7/(6z)\end{pmatrix}. \tag{3}
\]
The scalar equation for the first coordinate of the homogeneous system is
\[
 u''+\frac2z u'+\left(\frac5{36z^2}-4\right)u=0.
\]
The rational gauge $v=zu$ changes this to
$v''=(4-5/(36z^2))v$.

\begin{lemma}\label{lem:irred}
The coefficient differential module in (3) is irreducible over $\C(z)$.
\end{lemma}
\begin{proof}
The first coordinate is cyclic, since the off-diagonal entry is $-2$.
A factorization of its second-order scalar operator would give a rational
Riccati solution
\[
 w'+w^2=4-5/(36z^2).
\]
At a finite nonzero pole, comparison of the leading terms forces a simple
pole with residue $1$. At zero a pole is necessary, is simple, and its
residue $\rho$ satisfies $\rho(\rho-1)=-5/36$, so
$\rho=1/6$ or $5/6$. At infinity the polynomial part of $w$ is the
constant $2$ or $-2$. If $w=\pm2+b/z+O(z^{-2})$, the coefficient of $1/z$
in the Riccati equation gives $b=0$. But partial fractions give
$b=\rho+k$ for the number $k\ge0$ of other finite poles, a contradiction.
Thus the operator has no first-order factor; the rank-two module is simple.
\end{proof}

Let $\mathcal N_b$ be the $\C(z)$-span of all actual polynomial-amplitude
integrals on the three specified edges, together with their endpoint
exponentials. This span is finite dimensional and differential. On the
vertical edge the phase is $Y^2$; on the other edges, parametrized by $X$,
it is $g(X)=X^3+X^2/4-4X+1$. Ordinary one-variable division by
$\partial_s+zp'(s)$ reduces amplitudes to degree at most $\deg p-2$,
and an exact term gives the endpoint exponentials. The resulting closed
boundary system has coefficients in $\Q[z,z^{-1}]$. In particular there
are no finite poles except zero. Put $\mathcal N=\mathcal N_b+\C(z)$,
so the constant function $1$ is included.

The integration map from the coefficient module to meromorphic functions
modulo $\mathcal N_b$ is a differential map by (3). It is nonzero:
\[
 q+2=(X-1)^2(X+2)+Y^2
\]
shows a unique minimum $-2$ at the interior point $(1,0)$. Its minimum $d$
on the compact specified boundary satisfies $d>-2$. For large positive $s$,
every $B\in\mathcal N_b$ obeys $|B(-s)|\le Cs^N e^{-ds}$, whereas a small
positive-area sublevel neighborhood gives
$U_0(-s)\ge A_\epsilon e^{(2-\epsilon)s}$ for $0<\epsilon<d+2$.
Thus $U_0\notin\mathcal N_b$. Lemma~\ref{lem:irred} makes the map injective.
Its rank-two simple image intersects the differential module
$(\mathcal N_b+\C(z))/\mathcal N_b$, of rank at most one, trivially.
Hence the map stays injective modulo $\mathcal N$.

The interior coordinates, the reduced edge moments, their endpoint
exponentials and $1$ form a closed vector of entire $E$-functions over
$\Q[z,z^{-1}]$. To check the arithmetic condition directly, parametrize
$T$ by $X=2u$, $Y=(1-u)(2v-1)$ on $[0,1]^2$, with Jacobian $4(1-u)$.
Each Taylor coefficient before division by $m!$ is an algebraic integral
of a polynomial of degree $O(m)$. Common denominators through index $m$
divide $c^{m+1}\operatorname{lcm}(1,\ldots,Cm+C_0)^2$ for fixed integers,
and the conjugate sizes grow exponentially. The edges have the same
one-variable property. This verifies the $E$-function requirements.

At every nonzero algebraic $\xi$ the whole system is ordinary. Beukers'
\href{https://arxiv.org/html/math/0405549}{Theorem 3.2 and Lemma 3.1}
lift a numerical linear relation to a rational relation regular at $\xi$
with the prescribed specialization, also for dependent coordinate lists.
Every such functional relation has its two interior coefficients identically
zero by the preceding injectivity. We have therefore proved:

\begin{theorem}\label{thm:linear}
At every $\xi\in\Qb^\times$, $U_0(\xi),U_1(\xi)$ are linearly independent
over $\Qb$ modulo the algebraic span of the specified boundary values and
$1$. For every polynomial amplitude $r\in\Qb[X,Y]$,
\[
 (a_r(1/\xi),b_r(1/\xi))\ne(0,0)
 \quad\Longrightarrow\quad \int_T r e^{\xi q}\notin\Qb.
\]
A zero remainder only puts the integral in the boundary span. It makes
no zero-period assertion.
\end{theorem}

For example at $\xi=1$, $r=1+2Y^2$ has zero specialized remainder because
$r=(\partial_Y+2Y)Y$. Its integral is strictly positive. This is a
countercontrol to discarding the actual boundary term in (2).

\section{All three integer shape exponents at once}
The barycentric coordinates of $T$, in the order of vertices
$(0,-1),(0,1),(2,0)$, are
\[
 u=(2-X-2Y)/4,\qquad v=(2-X+2Y)/4,\qquad w=X/2.
\]
At the critical point $(1,0)$ they equal $(1/4,1/4,1/2)$, all units
modulo five. For nonnegative integers $a,b,c$, let $r_{abc}=u^av^bw^c$.
Its remainder lies in $\mathbb Z[1/6,h]^2$, and at $h=0$
\[
 a_{r_{abc}}(0)+b_{r_{abc}}(0)=r_{abc}(1,0)=4^{-a-b}2^{-c}. \tag{4}
\]
In particular it is a unit at five, independently of all three exponents.

\begin{theorem}\label{thm:all-shapes}
Let $\xi\ne0$ be algebraic. If $v_{\mathfrak p}(\xi)<0$ at some place
above five, then for every $a,b,c\ge0$,
\[
 \int_T u^av^bw^c e^{\xi q}\,dX\,dY\notin\Qb,
 \qquad\text{in particular it is nonzero.}
\]
\end{theorem}
\begin{proof}
At that place $h=1/\xi$ reduces to zero. All normal-form coefficients are
integral there. If both specialized remainder coordinates were zero,
reduction and (4) would contradict a nonzero residue. Apply
Theorem~\ref{thm:linear}. This is an all-exponent argument, not an
extrapolation from tests. It uses a strictly negative valuation; no claim
at every five-adic integral scale is inferred.
\end{proof}

Examples are $\xi=Ni/5$ for any integer $N$ not divisible by five, including
arbitrarily large imaginary scales. The executable certificate tests a
sufficient and necessary condition for the existence of such a place:
the monic minimal polynomial over $\Q$ has a coefficient nonintegral at
five. Indeed all conjugate roots are integral in $\overline{\Q}_5$ if and
only if its monic coefficients are integral. A failed test means only
that this arithmetic criterion is unavailable.

For a fixed triple $a,b,c$, its polynomial remainder vector is not
identically zero by (4). Thus the common-zero exceptional set
\[
 E_{abc}=\{\xi\in\C^\times:
          a_{r_{abc}}(1/\xi)=b_{r_{abc}}(1/\xi)=0\}
\]
is finite and algebraic. Every member is integral at every place above
five. Outside this set every algebraic scale has a nonzero period by
Theorem~\ref{thm:linear}. Emptiness of all $E_{abc}$ is not asserted.

\section{Factorial, Gamma and Gaussian consequences}
Let $Y_j$ be independent $\operatorname{Gamma}(\alpha_j,1)$ variables,
with $\alpha_j>0$, and partition their indices into three nonempty blocks.
Write their sums as $B_1,B_2,B_3$, their shape sums as $A_1,A_2,A_3$, and
$S=B_1+B_2+B_3$, $A=A_1+A_2+A_3$. Assume the \emph{three block shape sums}
$A_1,A_2,A_3$ are positive integers. In particular every positive integer
individual shape vector is allowed. Put
\[
 G=8B_3^3-6B_3S^2+(B_2-B_1)^2S.
\]
The angular variables $(B_1/S,B_2/S,B_3/S)$ have Dirichlet shapes
$(A_1,A_2,A_3)$, and
$G/S^3=q(2B_3/S,(B_2-B_1)/S)$.

\begin{corollary}\label{cor:family}
For $D\ge4$, an algebraic $\xi\ne0$ with a negative valuation at a place
above five, and any complex polynomial $R$ of degree at most $D-2$, set
\[
 F=S^D+D\xi S^{D-4}G+R.
\]
Then every sufficiently large $m$ satisfies $\mathbb E[F^m]\ne0$.
The same holds at $\xi=0$. This gives factorial families for every $n\ge3$,
in degree four already, for any partition into three nonempty blocks.
\end{corollary}
\begin{proof}
The integer-total-shape radial limit proved in the preceding note gives
\[
 \frac{\mathbb E[F^m]}{(A)_{Dm}}
 \longrightarrow \mathbb E_{\operatorname{Dir}(A_1,A_2,A_3)}e^{\xi q}
 =\frac{\Gamma(A)}{4\Gamma(A_1)\Gamma(A_2)\Gamma(A_3)}
 \int_Tu^{A_1-1}v^{A_2-1}w^{A_3-1}e^{\xi q}\,dX\,dY.
\]
The radial expansion only needs total shape $A\ge1$; its domination and
the compact angular measure still hold when individual shapes are not
integers. The angular grouping is exact. The limit is nonzero by
Theorem~\ref{thm:all-shapes}, independently of $R$. At zero the probability
limit is one. A sequence converging to a nonzero number is eventually nonzero.
\end{proof}

For example take $\xi=101i/5$, $D=4$, all $\alpha_j=1$, and
$B_1=Y_1,B_2=Y_2,B_3=Y_3+\cdots+Y_n$. Every polynomial
\[
 S^4+\frac{404i}{5}\{8B_3^3-6B_3S^2+(B_2-B_1)^2S\}+R_{\le2}
\]
has eventually nonzero factorial moments, in every dimension $n\ge3$.
The lower polynomial may involve every variable with arbitrary complex
coefficients. The angular Hessian has determinant $12X$, so this cubic
cannot be a polynomial in a single affine linear form.

For fixed block shapes the same eventual conclusion holds at every
algebraic $\xi\notin E_{A_1-1,A_2-1,A_3-1}$. If the individual shapes are
also algebraic (in particular positive integers), then at an arbitrary
complex scale outside that finite set, algebraic specialization of the polynomial moment
equations gives the FC conclusion. Explicitly, if all moments vanished,
invert $\xi$ and one nonzero specialized remainder coordinate and specialize
the finitely generated coefficient algebra to $\Qb$; Noetherianity reduces
the moment ideal to finitely many generators. At that algebraic point the
nonzero radial limit gives a contradiction. Positive powers retain the
same angular phase and scale:
\[
 F^k=S^{Dk}+Dk\xi S^{Dk-4}G+R',\qquad \deg R'\le Dk-2.
\]
Thus infinitely many moments survive at all complex scales outside the
finite exceptional set, including every transcendental scale. This last
conclusion does not assert a nonzero limit there. At an exceptional scale
the present proof makes no FC decision.

The Gaussian substitution $B_j=\frac12\sum_{\ell=1}^{2A_j}X_{j\ell}^2$
therefore gives the corresponding three-even-block Gaussian families.
The moment bridge preserves all original polynomial multipliers through
orthogonal averaging, as proved in the earlier transfer note; a nonzero
unit-period limit here is not a general theorem about arbitrary markers.

\section{Exact verification and remaining scope}
\href{two_angular.py}{The implementation} exports the actual polynomial
Stokes identity (1), an arithmetic certificate for all integer shape
triples, and fixed-amplitude specialized remainder certificates. Tests
check these identities, the rank-two connection and scalar equation,
the Riccati residues, the triangle/Dirichlet constants, grouped Gamma and
Gaussian moments, power closure and the zero-remainder boundary control.
Finite tests support the implementation; all-exponent nonvanishing uses
the written argument (4), irreducibility and arithmetic lifting.

Only the linear/extremum mechanism needed here has been audited from
EIT V. Its full independence theorem, continued-density elliptic detector,
rank-six example and all EIT VI block families remain outside this audit.
General two-angular-coordinate phases, arbitrary noninteger block shape
sums, unrestricted FC in dimension at least three, unrestricted SFC and
the complete bibliography remain open. No first surviving moment,
uniform consecutive window, or outside review is claimed.
\end{document}
